\documentclass[../../main.tex]{subfiles}
\begin{document}

\section*{问题}\label{section:ch03-问题}

\begin{exercise}\label{exercise:ch03-3-1}
设 $M$ 和 $N$ 是带边或无边光滑流形，$F:M\to N$ 是一个光滑映射。证明：对每个 $p\in M$，$\mathrm{d}F_{p}:T_{p}M\to T_{F(p)}N$ 是零映射，当且仅当 $F$ 在 $M$ 的每个连通分支上为常值。
\end{exercise}

\begin{exercise}\label{exercise:ch03-3-2}
证明命题~\cref{proposition:ch03-乘积流形切空间}（乘积流形的切空间）。
\end{exercise}

\begin{exercise}\label{exercise:ch03-3-3}
证明：若 $M$ 和 $N$ 是光滑流形，则 $T(M\times N)$ 微分同胚于 $TM\times TN$。
\end{exercise}

\begin{exercise}\label{exercise:ch03-3-4}
证明 $TS^{1}$ 微分同胚于 $S^{1}\times\mathbb{R}$。
\end{exercise}

\begin{exercise}\label{exercise:ch03-3-5}
令 $S^{1}\subset\mathbb{R}^{2}$ 为单位圆，令 $K\subset\mathbb{R}^{2}$ 为以原点为中心、边长为 $2$ 的正方形的边界：$K=\{(x,y):\max(|x|,|y|)=1\}$。证明存在一个同胚 $F:\mathbb{R}^{2}\to\mathbb{R}^{2}$ 使 $F(S^{1})=K$，但不存在具有同样性质的微分同胚。[提示：令 $\gamma$ 为一条像落在 $S^{1}$ 中的光滑曲线，并考虑 $\mathrm{d}F(\gamma'(t))$ 对坐标函数 $x$ 和 $y$ 的作用。]（用于第~123 页。）
\end{exercise}

\begin{exercise}\label{exercise:ch03-3-6}
在通常的 $\mathbb{C}^{2}\leftrightarrow\mathbb{R}^{4}$ 等同下，把 $S^{3}$ 看作 $\mathbb{C}^{2}$ 中的单位球面。对每个 $z=(z^{1},z^{2})\in S^{3}$，由 $\gamma_{z}(t)=(\mathrm{e}^{\mathrm{i}t}z^{1},\mathrm{e}^{\mathrm{i}t}z^{2})$ 定义一条曲线 $\gamma_{z}:\mathbb{R}\to S^{3}$。证明 $\gamma_{z}$ 是一条其速度永不为零的光滑曲线。
\end{exercise}

\begin{exercise}\label{exercise:ch03-3-7}
设 $M$ 是一个带边或无边光滑流形，$p$ 是 $M$ 的一点。令 $C^{\infty}_{p}(M)$ 记 $p$ 处光滑实值函数芽的代数，令 $D_{p}M$ 记 $C^{\infty}_{p}(M)$ 的导子所成向量空间。由 $(\Phi v)f=v([f]_{p})$ 定义一个映射 $\Phi:D_{p}M\to T_{p}M$。证明 $\Phi$ 是一个同构。（用于第~71 页。）
\end{exercise}

\begin{exercise}\label{exercise:ch03-3-8}
设 $M$ 是一个带边或无边光滑流形，$p\in M$。令 $V_{p}M$ 记从 $p$ 出发的光滑曲线在关系 $\gamma_{1}\sim\gamma_{2}$（即对每个定义在 $p$ 邻域上的光滑实值函数 $f$ 有 $(f\circ\gamma_{1})'(0)=(f\circ\gamma_{2})'(0)$）下的等价类所成集合。证明由 $\Psi[\gamma]=\gamma'(0)$ 定义的映射 $\Psi:V_{p}M\to T_{p}M$ 良定义且为双射。（用于第~72 页。）
\end{exercise}

\end{document}